\begin{document}

\Large{Who's affraid of subjective probability?. Maria Carla Galavotti. From Arithmetic to metaphysics.}

\begin{quote}
	n a “Technical report” on a sensitive subject, namely the chance of California being hit by an earthquake before 2030, statisticians David Freedman and Philip Stark criticize the subjective interpretation of probability – which they call ‘Bayesian’, following a currently widespread but disputable tendency because the two categories do not overlap since many Bayesians do not embrace the subjective theory. They write that “According to Bayesians, probability means degree of belief. This is measured on a scale running from 0 to 1. p.~151
\end{quote}

\section{del Artículo de KAhnemann y Tversky}

Thus, we suggest that mental simu-
lation yields a measure of the propensity of one's model of the situation to
generate various outcomes, much as the propensities of a statistical model
can be assessed by Monte Carlo techniques

We first list
some judgmental activities in which mental simulation appears to be
involved. We then describe a study of the cognitive rules that govern the
mental undoing of past events, and we briefly discuss the implications of
these rules for emotions that arise when reality is compared with a favored
alternative, which one had failed to reach but could easily imagine
reaching. We conclude this brief sketch of the simulation heuristic by
some remarks on scenarios, and on the biases that are likely to arise when
this heuristic is used.

The alterations that people introduce in stories can be classified as
downhill, uphill, or horizontal changes. A downhill change is one that
removes a surprising or unexpected aspect of the story, or otherwise
increases its internal coherence. An uphill change is one that introduces
unlikely occurrences. A horizontal change is one in which an arbitrary
value of a variable is replaced by another arbitrary value, which is neither
more nor less likely than the first. The experimental manipulation caused
a change of route to be downhill in one version, uphill in the other, with a
corresponding variation in the character of changes of the timing of Mr.
Jones's fatal trip. The manipulation was clearly successful: Subjects were
more likely to undo the accident by restoring a normal value of a variable
than by introducing an exception. In general, uphill changes are relatively
rare in the subjects' responses, and horizontal changes are non-existent



\section{DEl blog sobre Newton}

Esto es directo de la Principia. El énfasis está en el famoso hypothesis non fingo

Thus far I have explained the phenomena of the heavens and of our sea by the force of
gravity, but I have not yet assigned a cause to gravity. Indeed, this force arises
from some cause that penetrates as far as the centres of the sun iand planets without
any diminution of its power to act, and that acts not in proportion to the quantity
of the surfaces of the particles on which it acts (as mechanical causes are wont to
do) but in proportion to the quantity of solid matter, and whose action is extended
everywhere to immense distances, always decreasing as the squares of the distances.
Gravity toward the sun is compounded to the gravities towards the individual
particles of the sun, and at increasing distances from the sun decreases exactly as
the squares of the distances as far out as the orbit of Saturn, as is manifest from
the fact that the aphelia of the planets are at rest, and even as the far as the
farthest aphelia of the comets, providing that those aphelia are at rest. I have not
as yet been able to deduce from phenomena the reason for these properties, and I do
not feign hypotheses [ my emphasis]. For whatever is not deduced from the phenomena
must be called a hypothesis; and hypotheses, whether metaphysical or physical, or
based on occult qualities, or mechanical, have no place in experimental philosophy.
In this experimental philosophy, propositions are deduced from the phenomena and are
made general by induction. The impenetrability, mobility, and impetus of bodies, and
the laws of motion and the law of gravity have been found by this method. And it is
enough that gravity really exists and acts according to the laws that we have set
forth and is sufficient to explain all the motions of the heavenly bodies and of our
sea. 


\end{document}
